%% IEEE TAI edition — two-column working document.
%% Body is generated by build-tai.sh (pandoc -> body.tex) and \input below.
%% Double-anonymous: no author identity, affiliations, funding, or external links.
\documentclass[journal]{IEEEtran}

% --- xelatex font setup (system engine; Unicode §, ε, ∇ render natively) ---
\usepackage{fontspec}
\setmainfont{TeX Gyre Termes} % Times-metric-compatible -> realistic IEEE page count

\usepackage{amsmath,amssymb}
\usepackage{mathrsfs} % \mathscr
\usepackage[numbers,sort&compress]{natbib}
\usepackage{graphicx}
\usepackage{booktabs}
\usepackage{array}
\usepackage{multirow}
\usepackage{longtable}
\usepackage{calc} % pandoc table column widths use real-number factors
\providecommand{\real}[1]{#1}
\usepackage{url}
\usepackage{xcolor}
\usepackage[colorlinks=true,linkcolor=blue,citecolor=blue,urlcolor=blue]{hyperref}

% pandoc emits these helper macros but defines them only in its own template;
% with a fragment we must provide them ourselves.
\providecommand{\tightlist}{\setlength{\itemsep}{0pt}\setlength{\parskip}{0pt}}
\providecommand{\passthrough}[1]{#1}
\providecommand{\pandocbounded}[1]{#1}

% Keep manual Roman numerals in headings (# I. Introduction); disable auto-numbering.
\setcounter{secnumdepth}{0}

% Scale any oversized (pre-distillation) figures down to the column/page width
% so the baseline compiles without overfull boxes breaking layout.
\setkeys{Gin}{width=\linewidth,keepaspectratio}

\begin{document}

\title{Understanding Q-Learning and Deep Q-Learning in 2026:\\
A Methodological and Empirical Survey}

% ANONYMIZATION TOGGLE.
%   \anonfalse -> named working copy (authors visible; for internal circulation)
%   \anontrue  -> double-anonymous SUBMISSION build (REQUIRED for TAI; the
%                 Editorial Assistant returns non-anonymized submissions).
\newif\ifanon
\anontrue

\ifanon
  \author{Anonymous Submission to IEEE Transactions on Artificial Intelligence%
  \thanks{Double-anonymous review copy. Author identities, affiliations, and
  acknowledgements are withheld and provided via the submission portal.}}
  \markboth{IEEE Transactions on Artificial Intelligence,~Vol.~XX, No.~X, 2026}%
  {Understanding Q-Learning and Deep Q-Learning in 2026}
\else
  \author{Colby Wang, Ti, Divya, Kevin, Hamna, Logan,
  Eason Yishan Wu, Charles Jiahao Zhang, and Alp Guneysel%
  \thanks{E.~Y.~Wu is with Taipei, Taiwan Province of China.
  C.~J.~Zhang is with the University of Waterloo.
  Affiliations for the remaining authors are to be completed.
  \emph{Working copy --- set \texttt{\textbackslash anontrue} for the
  double-anonymous TAI submission build.}}}
  \markboth{IEEE Transactions on Artificial Intelligence,~Vol.~XX, No.~X, 2026}%
  {Wang \emph{et al.}: Understanding Q-Learning and Deep Q-Learning in 2026}
\fi

\maketitle

\begin{abstract}
Q-learning remains a cornerstone of Reinforcement Learning (RL), underpinning
many state-of-the-art deep RL algorithms. Yet, despite its foundational status,
no survey has systematically organized the diverse innovations that have shaped
Q-learning over the past three decades. This paper fills that gap by presenting
a comprehensive review of Q-learning variants, organized around the eight
foundational weaknesses of vanilla Q-learning that modern methods address ---
overestimation bias, sample inefficiency, brittle exploration, reward sparsity
and credit assignment, distribution shift, multi-agent coordination,
scaling and slow adaptation, and function-approximation instability. We introduce a unified
taxonomy covering both tabular and deep Q-learning methods, including offline
RL, multi-agent value decomposition, and distributed Q-learning, and compile
benchmark results from original papers across Atari and classic control
environments. To complement the literature review, we provide a comparative
analysis of six major open-source deep RL repositories --- Tianshou, XuanCe,
CleanRL, DQN Zoo, Stable Baselines3, and RLlib --- highlighting their coverage
of Q-learning algorithms and practical implementation challenges. Together,
these contributions offer a structured resource for researchers and
practitioners to navigate the evolution, empirical performance, and
implementation landscape of Q-learning.
\end{abstract}

\begin{IEEEkeywords}
Deep Learning, Machine Learning, Reinforcement Learning, Q-Learning, Survey
\end{IEEEkeywords}

\IEEEpeerreviewmaketitle

% --- Impact Statement (TAI requires 100--150 words after the abstract) ---
\section*{Impact Statement}
Reinforcement learning increasingly underpins decision-making systems in
robotics, autonomous vehicles, recommendation, and the post-training of large
language models --- and value-based Q-learning sits beneath much of this
technology. Yet practitioners face a fragmented literature: dozens of
algorithmic variants whose purpose, trade-offs, and software availability are
scattered across three decades of papers. This survey reorganizes that
landscape around the specific weaknesses each method was designed to fix,
letting an engineer move directly from a symptom --- overestimation, sparse
rewards, instability at scale --- to the methods that address it and the
open-source libraries that implement them. By pairing this diagnostic taxonomy
with calibrated benchmark evidence and a reproducibility audit of six widely
used codebases, the work lowers the barrier to selecting, deploying, and
extending Q-learning methods, and surfaces where research and tooling
investment are most needed.

% ===================== BODY (generated by build-tai.sh) =====================
\input{body.tex}

% ============================ BIBLIOGRAPHY ==================================
\bibliographystyle{IEEEtran}
\bibliography{refs}

\end{document}
